% ============================================================
% Section 5: Chinese Mathematics (Slides 32-36)
% ~200 BCE -- 1300 CE
% ============================================================
\section{Chinese Mathematics}

% ----------------------------------------------------------
% Slide 32: Section Divider
% ----------------------------------------------------------
{
\setbeamercolor{background canvas}{bg=chineseRed}
\begin{frame}[plain,shrink=22]
  \centering
  \vspace{1cm}
  \textcolor{white}{\Huge\bfseries China}\\[0.5em]
  \textcolor{white!90}{\Large Parallel Discoveries}\\[0.8em]
  \textcolor{white!75}{\large $\sim$200 BCE -- 1300 CE}\\[1.2em]

  % Chinese rod numerals illustration
  \visualplaceholder[5cm]{Horizontal rods for 1-5}

  \vspace{1em}
  % Timeline marker
  \visualplaceholder[5cm]{Diagram}

  \note{
    Section transition. Chinese mathematics developed largely independently of
    Greek and Indian traditions. Many results were discovered in China centuries
    before they appeared in Europe.
  }
\end{frame}
}

% ----------------------------------------------------------
% Slide 33: The Nine Chapters on the Mathematical Art
% ----------------------------------------------------------
\begin{frame}[shrink=22]{The Nine Chapters on the Mathematical Art}

  \begin{columns}[T]
    \begin{column}{0.52\textwidth}
      \begin{personbox}[\emph{Jiuzhang Suanshu}]{chineseRed}
        \small
        \textbf{Date:} Compiled c.\ 2nd century BCE -- 1st century CE\\
        \textbf{Origin:} Han Dynasty, China\\
        \textbf{Content:} 246 problems covering arithmetic, algebra, geometry, engineering
      \end{personbox}

      \vspace{0.4em}
      \begin{itemize}
        \item China's most important ancient mathematical text
        \item Includes \alert{Gaussian elimination} --- $\sim$2,000 years before Gauss!
        \item Methods for \textbf{square roots} and \textbf{cube roots}
        \item Practical problems: taxation, construction, surveying
        \item Nine chapters = nine topic areas
      \end{itemize}
    \end{column}

    \begin{column}{0.45\textwidth}
      \visualplaceholder[4.5cm]{Page from the \emph{Nine Chapters} showing mathematical problems (Wikimedia Commons / Chinese library digital collections)}

      \vspace{0.3em}
      % The "broken bamboo" problem
      \visualplaceholder[5cm]{Ground line}
    \end{column}
  \end{columns}

  \note{
    The Nine Chapters is comparable to Euclid's Elements in importance, but very
    different in style: practical problems and algorithms rather than axiomatic proofs.
    \verifytag{The dating is debated. The text was compiled over centuries; the final
    form is usually dated to the 1st century CE (Han Dynasty).}
    ``Gaussian elimination'' is so named because Gauss used it in the 1800s ---
    but it was in this book 2,000 years earlier.
  }
\end{frame}

% ----------------------------------------------------------
% Slide 34: Liu Hui
% ----------------------------------------------------------
\begin{frame}[shrink=17]{Liu Hui}

  \begin{columns}[T]
    \begin{column}{0.48\textwidth}
      \begin{personbox}[Liu Hui]{chineseRed}
        \small
        \textbf{Dates:} fl.\ c.\ 263 CE\\
        \textbf{Origin:} Cao Wei kingdom, China (Three Kingdoms period)\\
        \textbf{Contribution:} Commentary on the \emph{Nine Chapters} (263 CE);
        polygon method for $\pi$; principles of cartographic surveying
      \end{personbox}

      \vspace{0.4em}
      \begin{itemize}
        \item Used a \alert{polygon method} (like Archimedes) to calculate $\pi$
        \item Achieved $\pi \approx 3.14159$ (5 decimal places)
        \item Developed independently of Archimedes
        \item Also pioneered \textbf{map-making techniques}
      \end{itemize}
    \end{column}

    \begin{column}{0.49\textwidth}
      % Liu Hui's polygon method and comparison with Archimedes
      \visualplaceholder[5cm]{Circle}

      \vspace{0.3em}
      % Comparison table
      {\small
      \begin{tabular}{lcc}
        \toprule
        \textbf{Who} & \textbf{When} & \textbf{$\pi$ accuracy} \\
        \midrule
        Archimedes & c.\ 250 BCE & $3.1408\ldots$ \\
        Liu Hui & 263 CE & $3.14159$ \\
        Zu Chongzhi & c.\ 480 CE & $3.1415926\ldots$ \\
        \bottomrule
      \end{tabular}
      }
    \end{column}
  \end{columns}

  \note{
    Liu Hui's method is strikingly similar to Archimedes's --- but developed
    independently, 500 years later and 8,000 km away. Same idea, different civilisation.
    \verifytag{The exact decimal accuracy Liu Hui achieved is sometimes stated as
    3.1416 (4 places) or 3.14159 (5 places). Sources vary.}
  }
\end{frame}

% ----------------------------------------------------------
% Slide 35: Zu Chongzhi
% ----------------------------------------------------------
\begin{frame}[shrink=15]{Zu Chongzhi --- The World's Best $\pi$ for a Millennium}

  \begin{columns}[T]
    \begin{column}{0.52\textwidth}
      \begin{personbox}[Zu Chongzhi]{chineseRed}
        \small
        \textbf{Dates:} 429 -- 500 CE\\
        \textbf{Origin:} Liu Song / Southern Qi dynasties, China\\
        \textbf{Contribution:} Calculated $\pi$ to 7 decimal places ---
        the most accurate value in the world for nearly 1,000 years
      \end{personbox}

      \vspace{0.4em}
      \begin{itemize}
        \item $3.1415926 < \pi < 3.1415927$
        \item The fraction $\dfrac{355}{113}$ is accurate to \textbf{6 decimal places}
        \item This approximation was not bettered for \alert{$\sim$1,000 years}
        \item Extended Liu Hui's polygon method
      \end{itemize}
    \end{column}

    \begin{column}{0.45\textwidth}
      % The fraction 355/113 highlighted
      \visualplaceholder[5cm]{Diagram}

      \vspace{0.5em}
      % Timeline of pi accuracy
      \visualplaceholder[5cm]{Entries}
    \end{column}
  \end{columns}

  \note{
    ``His approximation 355/113 is so good that it was not bettered for a millennium.''
    Ask students: how do you even compute $\pi$ to 7 decimal places without
    a calculator? Answer: polygons with thousands of sides and enormous patience.
    This is also a good moment to note that the $\pi$ race was a genuinely
    global competition, with contributions from Egypt, Greece, India, and China.
  }
\end{frame}

% ----------------------------------------------------------
% Slide 36: Chinese Mathematics Legacy [EXTENDED]
% ----------------------------------------------------------
\begin{frame}[shrink=25]{Chinese Mathematics --- Legacy}

  \begin{columns}[T]
    \begin{column}{0.47\textwidth}
      \textbf{Independent discoveries:}
      \begin{itemize}
        \item \alert{Gaussian elimination} ($\sim$2,000 yrs before Gauss)
        \item \textbf{Horner's method} for polynomials
        \item \textbf{Yang Hui's triangle} (= Pascal's triangle)
        \item \textbf{Chinese Remainder Theorem} --- still used in modern cryptography
      \end{itemize}

      \vspace{0.3em}
      \textbf{Key figures:}\\[0.2em]
      {\small
      \textbf{Sunzi} (fl.\ c.\ 3rd--5th c.\ CE) --- Chinese Remainder Theorem\\
      \textbf{Yang Hui} (c.\ 1238--1298) --- described the triangle in 1261\\
      \textbf{Zhu Shijie} (c.\ 1249--1314) --- polynomial algebra with 4 unknowns
      }
    \end{column}

    \begin{column}{0.50\textwidth}
      % Yang Hui's Triangle (Pascal's Triangle)
      \visualplaceholder[5cm]{Triangle rows}

      \vspace{0.3em}
      {\tiny\textcolor{gray}{Yang Hui published this 392 years before Blaise Pascal.}}
    \end{column}
  \end{columns}

  \vspace{0.3em}
  \begin{center}
    \textcolor{islamicGreen}{\emph{``Now we arrive at one of the most transformative periods in
    mathematical history --- one that is too often overlooked in Western textbooks\ldots''}}
  \end{center}

  \note{
    Emphasise the pattern of independent discovery: Gaussian elimination, Pascal's triangle,
    and the Remainder Theorem all carry European names but were known in China centuries earlier.
    \verifytag{Sunzi dates: fl.\ c.\ 3rd--5th century CE is standard but uncertain.}
    \verifytag{Yang Hui dates: c.\ 1238--1298 CE --- verify.}
    \verifytag{Zhu Shijie dates: c.\ 1249--1314 CE --- verify.}
    Transition: ``Now we arrive at one of the most transformative periods in mathematical
    history --- the Islamic Golden Age\ldots''

    CRITICAL TRANSITION: After the final sentence, PAUSE. Then ask: ``Mathematics
    was flourishing in India and China for centuries. Every one of these results
    is real, verified, documented. But where is this story in your textbook?''
    [Beat of silence.] ``The next section will show you what else was hidden.''
  }
\end{frame}
